\section{Lean Formalization}
\label{app:lean}
\lstdefinelanguage{lean}{
  morekeywords={def,inductive,mutual,end,match,with,fun,forall,exists,theorem,structure,where,let,if,then,else,do,pure,abbrev,namespace,import,deriving,by,sorry},
  morecomment=[l]{--},
  sensitive=true
}
\lstset{breaklines=true}
\newcommand{\leanfile}[1]{\texttt{lean/Funnelling/#1}}

This appendix relates the formalism of Section~\ref{sec:core} and Section~\ref{sec:java} to a Lean~4 formalization in the folder \texttt{lean} of the paper repository.
The Lean formalization is set up, but the proofs of the soundness theorems are not done yet.

\subsection{Request}
This appendix answers the following request, reported verbatim, so that we can check later how much of it is addressed.
\begin{lstlisting}[language={},basicstyle=\small\ttfamily,breaklines=true,columns=fullflexible]

The Appendix, for each piece of formalism we report the following:

(A)Original formalism

(B)Corresponding Lean formalism

(C)English explanation of the differences and either
- why it ok for those differences to be there.
- or what the more appropriate form of the original formalism should be

(D)A clear report of the lean statements and an high level description of the proofs.

(E)A list of validation examples. Those are:
- programs that Do pass the lean type system AND the paper type system, or  that do not pass both the lean type system and the paper type system.
- programs that reduce in some steps equally in both the paper type system and the lean type system. Basically, this is 'testing': we are empirically checking that the two formalism are the same by checking that they behave the same on a range of examples including interesting corner cases.
\end{lstlisting}

\subsection{Status}
The Lean version is 4.35.0, without Mathlib; \texttt{lean/lean-toolchain} fixes it.
\texttt{lake build} in the folder \texttt{lean} checks all the files.
The files are:
\begin{itemize}
\item \leanfile{Syntax.lean}: types, funnels, expressions, methods, declarations, class declarations $K$.
\item \leanfile{Substitution.lean}: substitution of type variables and of variables, and the renaming used by $\mathsf{fresh}$.
\item \leanfile{Declarations.lean}: $\mathsf{declared}$, $\mathsf{scope}$, the pairing, $\mathrm{generic}$, $M\in C$, $\mathsf{typeOf}$, values, $\mathsf{fresh}$.
\item \leanfile{WellFormed.lean}: the well-formedness requirements A1-A11.
\item \leanfile{Typing.lean}: subtyping, kinding, checked casts, $\mathsf{overrideOk}$, $\mathsf{consistent}$, typing, method, class and program typing.
\item \leanfile{Reduction.lean}: evaluation contexts and the reduction.
\item \leanfile{Soundness.lean}: the soundness statements.
\item \leanfile{Examples.lean}: the validation examples.
\end{itemize}
Lean accepts all the files.
The three theorems \texttt{stepWellFormed}, \texttt{progress} and \texttt{preservation} are stated with \texttt{sorry}.
The corollary \texttt{typeSoundness} is proved from those three theorems.
Section~\ref{app:lean:examples} lists which validation examples are proved and which are stated with \texttt{sorry}.
The mini proof that checked the installation is not in the repository:
it proved $a+b=b+a$ with \texttt{omega}, $\mathsf{Even}\ 4$ for an inductive predicate $\mathsf{Even}$, $\neg\mathsf{Even}\ 1$ by \texttt{cases}, and evaluated an expression with \texttt{\#eval}.
\\[1ex]
The Lean code is written in ASCII:
\texttt{->} is $\rightarrow$, \texttt{/\char`\\} is $\wedge$, \texttt{\char`\\/} is $\vee$, \texttt{List.Mem x xs} is $x\in xs$, and \texttt{delta}, \texttt{gamma}, \texttt{sigma} are $\Delta$, $\Gamma$, $\sigma$.

\subsection{Differences Common to All the Pieces}
\label{app:lean:common}
The following differences hold for every piece below, and the pieces do not repeat them.
\begin{itemize}
\item \textbf{Names.}
Lean has one type \texttt{Name} (strings) for class names, method names, field names, variables and type variables.
The paper uses a different metavariable for each.
This is fine: no definition compares two names of different kinds.
\item \textbf{Grammar and well-formedness.}
In the paper, the metavariables range over the terms of the grammar, and the well-formedness requirements restrict the grammar (Section~\ref{sec:clarifications}).
In particular, a metavariable $v$ ranges over values and a metavariable $G$ ranges over ground types.
Lean defines the raw syntax as inductive types, and the restrictions as Boolean functions:
\texttt{Expr.isValue} for values, \texttt{Ty.isGround} for ground types, \texttt{wfDecls} and \texttt{wfExpr} for A1-A11.
Thus every Lean rule with a metavariable $v$ or $G$ of the paper has an explicit premise \texttt{isValue} or \texttt{allGround}.
The theorems assume \texttt{wfDecls} and \texttt{wfExpr} for the initial program.
The Lean reduction does not require the result to be well formed; the theorem \texttt{stepWellFormed} states that it is.
In the paper an instance of a reduction rule exists only when its result is an expression, so the paper needs no such lemma.
The two formulations agree when \texttt{stepWellFormed} holds.
\item \textbf{Sequences.}
The paper applies a judgement to sequences position by position, and requires the same length (Section~\ref{sec:declarations}).
Lean writes a list premise as a length equality and a premise for every pair in \texttt{List.zip}, or as a premise for every element of a list.
\item \textbf{Substitution and lengths.}
The paper defines $\tsubst{T}{\Delta}{\pl{T}}$ only when $\Delta$ and $\pl{T}$ have the same length.
Lean builds a substitution with \texttt{mkSubst xs ts}, which is \texttt{none} when the lengths differ; a premise \texttt{mkSubst xs ts = some sigma} requires the same length.
\item \textbf{Functions instead of relations.}
The paper defines $\mathsf{declared}$, $\mathsf{scope}$, $\mathrm{generic}$, the pairing and $\mathsf{typeOf}$ with \textbf{Define} boxes; those definitions have at most one result.
Lean defines them as functions, so the validation examples can compute them.
The inference rules (reduction, typing, kinding, subtyping, checked cast, method typing) are inductive relations in Lean, as in the paper.
\item \textbf{Equality of types.}
Lean derives no decidable equality for the nested inductive types; \texttt{Ty.beq} is a Boolean equality on types, used only by the pairing.
The proofs need the lemma \texttt{Ty.beq t t' = true <-> t = t'}.
\end{itemize}

\subsection{Syntax}
\paragraph*{Original.}
\[
\begin{array}{rcl}
e & ::= & x \mid \call{e}{\pl{T}}{m}{\pl{e}} \mid \explicitcall{e}{T}{\pl{T}}{m}{\pl{e}} \mid e\fstop f\\
&& \mid \obj{C}{\pl{T}}{\pl{e}}\\
&& \mid \innerobj{C}{\pl{F}}{C'}{\pl{T}'}{\pl{e}}{\pl{M}}\\
&& \mid \ifthen{e_0}{T}{x}{e_1}{e_2}\\
T & ::= & \gen{C}{\pl{T}} \mid X \qquad F ::= G \mid \funnel{X} \qquad G ::= \gen{C}{\pl{G}}\\
\Gamma & ::= & T_1\gap x_1,\ldots,T_n\gap x_n \qquad \Delta ::= \ext{X_1}{T_1},\ldots,\ext{X_n}{T_n}\\
\mathsf{Dec} & ::= & \subclass{C}{\Delta}{C'}{\pl{T}}{\pl{M}}{\pl{S}} \qquad S ::= T\gap f\semi\\
M & ::= & \method{\Delta}{T_0}{m}{T_1\gap x_1,\ldots,T_n\gap x_n}{\ret{e}}\\
K & ::= & \mathsf{Dec} \mid \innerobj{C}{\pl{F}}{C'}{\pl{T}}{\pl{e}}{\pl{M}}
\end{array}
\]
\paragraph*{Lean.}
\begin{lstlisting}[language=lean,basicstyle=\scriptsize\ttfamily]
inductive Ty where
  | cls (c : Name) (args : List Ty)
  | var (x : Name)
inductive Funnel where
  | ground (g : Ty)
  | funnel (x : Name)
abbrev Bounds := List (Prod Name Ty)
mutual
inductive Expr where
  | var (x : Name)
  | call (recv : Expr) (targs : List Ty) (m : Name) (args : List Expr)
  | superCall (recv : Expr) (start : Ty) (targs : List Ty) (m : Name) (args : List Expr)
  | field (recv : Expr) (f : Name)
  | new (c : Name) (targs : List Ty) (args : List Expr)
  | newLocal (c : Name) (fs : List Funnel) (sup : Name) (supArgs : List Ty)
      (args : List Expr) (ms : List Method)
  | ifInst (e0 : Expr) (t : Ty) (x : Name) (e1 e2 : Expr)
inductive Method where
  | mk (bounds : Bounds) (ret : Ty) (m : Name) (params : List (Prod Ty Name)) (body : Expr)
end
structure Dec where
  name : Name
  bounds : Bounds
  sup : Name
  supArgs : List Ty
  methods : List Method
  fields : List (Prod Ty Name)
inductive K where
  | top (dec : Dec)
  | loc (c : Name) (fs : List Funnel) (sup : Name) (supArgs : List Ty)
      (args : List Expr) (ms : List Method)
abbrev Decls := List Dec
abbrev Ctx := List (Prod Name Ty)
\end{lstlisting}
\paragraph*{Differences.}
Ground types are not a separate type: \texttt{Funnel.ground} takes any \texttt{Ty}, and \texttt{wfDecls} and \texttt{wfExpr} check that it is ground (Section~\ref{app:lean:common}).
A type environment $\Gamma$ is a list of pairs \texttt{(x, T)}, in the order variable then type; \texttt{gamma.lookup x} is $\Gamma(x)$.
Extending $\Gamma$ with $T\gap x$ appends \texttt{(x, T)}; \texttt{lookup} returns the first binding, and by A7 there is only one, except for \Q@this@, which the rules remove before adding a new one.
This is fine.

\subsection{Well-Formedness Requirements A1-A11}
\paragraph*{Original.}
A1-A11 of Section~\ref{sec:syntax}.
\paragraph*{Lean.}
\begin{lstlisting}[language=lean,basicstyle=\scriptsize\ttfamily]
def a1 (d : Decls) : Bool := nodup ((declaredKs d).map K.name)
def a2 (d : Decls) : Bool :=
  d.all fun dec => nodup ((superChain d (d.length + 1) dec.name).flatMap fun c => c.fields.map Prod.snd)
def a3 (ks : List K) : Bool := ks.all fun k => nodup (k.methods.map Method.name)
def a4 (ms : List Method) : Bool := ms.all fun m => nodup m.params
def a5 (d : Decls) (ms : List Method) : Bool :=
  d.all (fun dec => nodup (dom dec.bounds)) && ms.all fun m => nodup (dom m.bounds)
def a6 (ks : List K) : Bool := ks.all fun k => nodup (domF k.funnels)
def a7Decls (d : Decls) : Bool := d.all fun dec => Method.noShadowList (dom dec.bounds) [] dec.methods
def a8 (ms : List Method) (ifVars : List Name) : Bool :=
  ms.all (fun m => !(m.params.contains thisName)) && !(ifVars.contains thisName)
def a9 (d : Decls) : Bool :=
  (d.any fun dec => dec.name == objectName && dec.bounds.isEmpty && dec.sup == objectName &&
    dec.supArgs.isEmpty && dec.methods.isEmpty && dec.fields.isEmpty) &&
  d.all fun dec => reachesObject d dec.name
def a10 (d : Decls) : Bool := (declaredKs d).all fun k => (findDec d k.sup).isSome
def a11 (d : Decls) : Bool :=
  (declaredKs d).all fun k => k.funnels.all fun
    | .funnel _ => true
    | .ground _ => false
def groundFunnels (ks : List K) : Bool :=
  ks.all fun k => k.funnels.all fun
    | .funnel _ => true
    | .ground g => g.isGround
def wfDecls (d : Decls) : Bool :=
  a1 d && a2 d && a3 (declaredKs d) && a4 d.methods && a5 d d.methods && a6 (declaredKs d) &&
  a7Decls d && a8 d.methods (d.flatMap fun dec => Method.ifVarsList dec.methods) &&
  a9 d && a10 d && a11 d && groundFunnels (declaredKs d)
def wfExpr (e : Expr) : Bool :=
  a3 e.locals && a4 e.methods && a5 [] e.methods && a6 e.locals && e.noShadow [] [] &&
  a8 e.methods e.ifVars && groundFunnels e.locals
\end{lstlisting}
\texttt{Expr.noShadow tvs vs e} (A7) walks $e$ with the type variables \texttt{tvs} and the variables \texttt{vs} in scope.
It fails when a method declares a type variable in \texttt{tvs} or a parameter in \texttt{vs}, or when an \Q@if@ binds a variable in \texttt{vs}.
It removes \Q@this@ from \texttt{vs} when it enters the methods of a local class.
\texttt{superChain d n c} is the chain of superclasses of $c$, cut after $n$ classes.
\paragraph*{Differences.}
\texttt{wfDecls} checks the requirements on \decls; \texttt{wfExpr} checks the requirements that also restrict the main expression (A3-A8).
A2 and A9 follow the chain of superclasses at most $|\decls|+1$ times: an acyclic chain is never longer, and A9 rejects a cyclic chain, since a cyclic chain does not reach \Q@Object@.
\texttt{groundFunnels} is the grammar of $G$, not a requirement of the paper (Section~\ref{app:lean:common}).
This is fine.

\subsection{Values and Evaluation Contexts}
\paragraph*{Original.}
\[
\begin{array}{rcl}
v & ::= & \obj{C}{\pl{G}}{\pl{v}} \mid \innerobj{C}{\pl{G}}{C'}{\pl{G}'}{\pl{v}}{\pl{M}}\\
\mathcal{E} & ::= & \square \mid \call{\mathcal{E}}{\pl{G}}{m}{\pl{e}} \mid \call{v}{\pl{G}}{m}{\pl{v},\mathcal{E},\pl{e}}\\
&& \mid \explicitcall{\mathcal{E}}{G}{\pl{G}}{m}{\pl{e}} \mid \explicitcall{v}{G}{\pl{G}}{m}{\pl{v},\mathcal{E},\pl{e}}\\
&& \mid \mathcal{E}\fstop f \mid \obj{C}{\pl{G}}{\pl{v},\mathcal{E},\pl{e}}\\
&& \mid \innerobj{C}{\pl{G}}{C'}{\pl{G}'}{\pl{v},\mathcal{E},\pl{e}}{\pl{M}}\\
&& \mid \ifthen{\mathcal{E}}{G}{x}{e_1}{e_2}
\end{array}
\]
\paragraph*{Lean.}
\begin{lstlisting}[language=lean,basicstyle=\scriptsize\ttfamily]
mutual
def Expr.isValue : Expr -> Bool
  | .new _ gs vs => Ty.allGround gs && Expr.allValues vs
  | .newLocal _ fs _ gs vs _ => fs.all Funnel.isGround && Ty.allGround gs && Expr.allValues vs
  | _ => false
def Expr.allValues : List Expr -> Bool
  | [] => true
  | e :: es => e.isValue && Expr.allValues es
end
inductive ECtx where
  | hole
  | callRecv (ctx : ECtx) (gs : List Ty) (m : Name) (es : List Expr)
  | callArg (v : Expr) (gs : List Ty) (m : Name) (vs : List Expr) (ctx : ECtx) (es : List Expr)
  | superRecv (ctx : ECtx) (g : Ty) (gs : List Ty) (m : Name) (es : List Expr)
  | superArg (v : Expr) (g : Ty) (gs : List Ty) (m : Name) (vs : List Expr) (ctx : ECtx) (es : List Expr)
  | field (ctx : ECtx) (f : Name)
  | newArg (c : Name) (gs : List Ty) (vs : List Expr) (ctx : ECtx) (es : List Expr)
  | localArg (c : Name) (fs : List Funnel) (c' : Name) (gs : List Ty) (vs : List Expr) (ctx : ECtx)
      (es : List Expr) (ms : List Method)
  | ifCond (ctx : ECtx) (g : Ty) (x : Name) (e1 e2 : Expr)
\end{lstlisting}
\texttt{ECtx.plug E e} is $\mathcal{E}[e]$, and \texttt{ECtx.ok E} checks that the $v$ are values and the $G$ are ground.
\paragraph*{Differences.}
None beyond Section~\ref{app:lean:common}.

\subsection{Declared Classes, Scope, Pairing and Generic}
\paragraph*{Original.}
The definitions of $\declared{K}$, $\mathsf{scope}(\gen{C}{\pl{T}})$, $\pairing{\Delta}{\pl{F}}{\pl{T}}$, $M\in K$, $M\in C$, $\generic{C}$ and $|F|$ in Section~\ref{sec:declarations}.
\paragraph*{Lean.}
\begin{lstlisting}[language=lean,basicstyle=\scriptsize\ttfamily]
def pairing (delta : Bounds) : List Funnel -> List Ty -> Option (Prod Bounds (List Ty))
  | [], [] => some ([], [])
  | .funnel x :: fs, t :: ts => do
    let b <- delta.lookup x
    let (delta', ts') <- pairing delta fs ts
    pure ((x, b) :: delta', t :: ts')
  | .ground g :: fs, t :: ts => if g.beq t then pairing delta fs ts else none
  | _, _ => none
def declaredKs (d : Decls) : List K :=
  d.flatMap fun dec => K.top dec :: Method.localsList dec.methods
def Declared (d : Decls) (k : K) : Prop := List.Mem k (declaredKs d)
def scope (d : Decls) (c : Name) (ts : List Ty) : Option (Prod Bounds (List Ty)) :=
  match findDec d c with
  | some dec => some (dec.bounds, ts)
  | none =>
    match (d.flatMap fun dec => Method.scopeEntriesList dec.bounds dec.methods).lookup c with
    | some (env, fs) => pairing env fs ts
    | none => none
def generic (d : Decls) (c : Name) : Option (List Name) :=
  match findK d c with
  | some (.top dec) => some (dom dec.bounds)
  | some (.loc _ fs _ _ _ _) => some (domF fs)
  | none => none
def MethodIn (d : Decls) (mt : Method) (c : Name) : Prop :=
  exists k, Declared d k /\ k.name = c /\ List.Mem mt k.methods
\end{lstlisting}
\texttt{Expr.locals e} lists every local class written anywhere in $e$, also inside the methods of local classes.
\texttt{Expr.scopeEntries env e} lists, for every local class $C$ written in $e$, the name $C$, the bounds \texttt{env} in scope where $C$ is written, and the funnels of $C$.
In the methods of $C$, the bounds in scope are the bounds of the pairing of \texttt{env} with the funnels of $C$, followed by the bounds of the method.
\paragraph*{Differences.}
\texttt{declaredKs} lists every local class written anywhere in the methods of \decls.
The paper reaches the same local classes through a chain of $\mathsf{declared}$, where each local class occurs in a method of the previous one.
\texttt{scope} computes, for a local class, the same bounds as the paper: the bounds of $\mathsf{scope}(\generic{C})$ for the enclosing class $C$, followed by the bounds of the method.
By A1 there is one entry per class name, so \texttt{lookup} finds the only one.
\texttt{generic} returns the type variables $\pl{X}$ of $\generic{C}=\gen{C}{\pl{X}}$; for a local class, $\mathrm{dom}(\pl{F})$ is $|\pl{F}|$ by A11.
This is fine.

\subsection{Substitution of Type Variables and of Variables}
\paragraph*{Original.}
The definitions of Sections~\ref{sec:substitution} and~\ref{sec:termsubstitution}.
\paragraph*{Lean.}
\begin{lstlisting}[language=lean,basicstyle=\scriptsize\ttfamily]
def mkSubst (xs : List Name) (ts : List Ty) : Option TSubst :=
  if xs.length = ts.length then some (xs.zip ts) else none
def mkSubstF (fs : List Funnel) (ts : List Ty) : Option TSubst :=
  if fs.length = ts.length then
    some ((fs.zip ts).filterMap fun
      | (.funnel x, t) => some (x, t)
      | (.ground _, _) => none)
  else none
def Ty.subst (sigma : TSubst) : Ty -> Ty
  | .var x => match sigma.lookup x with | some t => t | none => .var x
  | .cls c ts => .cls c (Ty.substList sigma ts)
def Funnel.subst (sigma : TSubst) : Funnel -> Funnel
  | .ground g => .ground (g.subst sigma)
  | .funnel x => match sigma.lookup x with | some t => .ground t | none => .funnel x
def Expr.tsubst (sigma : TSubst) : Expr -> Expr
  ...
  | .newLocal c fs c' ts es ms =>
    .newLocal c (fs.map (Funnel.subst sigma)) c' (Ty.substList sigma ts)
      (Expr.tsubstList sigma es) (Method.tsubstList sigma ms)
def Expr.vsubst (sigma : VSubst) : Expr -> Expr
  | .var x => match sigma.lookup x with | some v => v | none => .var x
  ...
  | .newLocal c fs c' ts es ms =>
    .newLocal c fs c' ts (Expr.vsubstList sigma es)
      (Method.vsubstList (sigma.filter fun p => p.1 != thisName) ms)
\end{lstlisting}
The cases written as \texttt{...} go into every subexpression, as in the paper.
\paragraph*{Differences.}
A substitution is a list of pairs; \texttt{lookup} returns the first pair, and \texttt{mkSubst} builds a substitution with distinct names when the names are distinct (A5, A6).
\texttt{Funnel.subst} turns \texttt{X funnel} into \texttt{.ground T} without checking that $T$ is ground; the paper substitutes only ground types into expressions, and \texttt{stepWellFormed} states that the results are ground.
The variable substitution into the methods of a local class drops \Q@this@, as in the paper.
This is fine.

\subsection{Renaming: $\mathsf{fresh}$}
\paragraph*{Original.}
$\mathsf{fresh}(e,e')$ holds when $e'$ is $e$ where every name bound in $e$ is renamed into a fresh name, that is a name occurring nowhere in the program; distinct binders get distinct fresh names; \Q@this@ is never renamed.
\paragraph*{Lean.}
\begin{lstlisting}[language=lean,basicstyle=\scriptsize\ttfamily]
def freshNames (avoid ns : List Name) : Bool :=
  decide ns.Nodup && ns.all fun n => !(avoid.contains n)
def Fresh (avoid : List Name) (e e' : Expr) : Prop :=
  exists ns, Expr.freshen [] [] ns e = some (e', []) /\ freshNames avoid ns = true
def FreshM (avoid : List Name) (mt mt' : Method) : Prop :=
  exists ns, Method.freshen [] [] ns mt = some (mt', []) /\ freshNames avoid ns = true
def avoidNames (d : Decls) (e : Expr) : List Name := d.names ++ e.names
\end{lstlisting}
\texttt{Expr.freshen rv rt ns e} walks $e$, takes the next name of \texttt{ns} for every binder, and renames the bound occurrences: \texttt{rv} renames variables and \texttt{rt} renames type variables.
It returns the renamed expression and the unused names.
\paragraph*{Differences.}
The paper avoids the names of the whole program.
The Lean rules avoid the names of \decls and of the redex (\texttt{avoidNames}).
Fewer names to avoid is enough for A7: the hole of an evaluation context is under no binder, so a name of the main expression outside the redex never binds around the result.
This is fine, and the paper text can stay as it is.

\subsection{Reduction}
\paragraph*{Original.}
The rules of Section~\ref{sec:opsem}.
\paragraph*{Lean.}
\begin{lstlisting}[language=lean,basicstyle=\scriptsize\ttfamily]
inductive Step (d : Decls) : Expr -> Expr -> Prop
  | simpleFieldFound {c gs vs1 vs f v} {dec : Dec} :
      List.Mem dec d -> dec.name = c ->
      vs1.length = dec.fields.length ->
      (Expr.new c gs (vs1 ++ vs)).isValue = true ->
      ((fieldNames dec).zip vs1).lookup f = some v ->
      Step d (.field (.new c gs (vs1 ++ vs)) f) v
  | simpleFieldSuper {c gs vs1 vs f sigma} {dec : Dec} :
      List.Mem dec d -> dec.name = c ->
      vs1.length = dec.fields.length ->
      (Expr.new c gs (vs1 ++ vs)).isValue = true ->
      (fieldNames dec).contains f = false ->
      Not (c = objectName) ->
      mkSubst (dom dec.bounds) gs = some sigma ->
      Step d (.field (.new c gs (vs1 ++ vs)) f) (.field (.new dec.sup (Ty.substList sigma dec.supArgs) vs) f)
  | innerFieldSuper {c fs c' gs vs ms f} :
      (Expr.newLocal c fs c' gs vs ms).isValue = true ->
      Step d (.field (.newLocal c fs c' gs vs ms) f) (.field (.new c' gs vs) f)
  | ifThen {v c' ts ts' tv x e1 e1' e2} :
      v.isValue = true -> typeOf v = some tv ->
      Sub d [] tv (.cls c' ts') ->
      Fresh (avoidNames d (.ifInst v (.cls c' ts) x e1 e2)) e1 e1' ->
      Step d (.ifInst v (.cls c' ts) x e1 e2) (e1'.vsubst [(x, v)])
  | ifElse {v c' ts tv x e1 e2} :
      v.isValue = true -> typeOf v = some tv ->
      (forall ts', Not (Sub d [] tv (.cls c' ts'))) ->
      Step d (.ifInst v (.cls c' ts) x e1 e2) e2
  | simpleAddSuper {c gs vs0 gs' m vs} :
      (Expr.new c gs vs0).isValue = true -> Ty.allGround gs' = true -> Expr.allValues vs = true ->
      Step d (.call (.new c gs vs0) gs' m vs) (.superCall (.new c gs vs0) (.cls c gs) gs' m vs)
  | simpleFound {v0 c gs gs' m vs mt delta' t0 ps e0 sigma} {dec : Dec} :
      List.Mem dec d -> dec.name = c ->
      v0.isValue = true -> Ty.allGround gs = true -> Ty.allGround gs' = true -> Expr.allValues vs = true ->
      List.Mem mt dec.methods ->
      FreshM (avoidNames d (.superCall v0 (.cls c gs) gs' m vs)) mt (.mk delta' t0 m ps e0) ->
      mkSubst (dom dec.bounds ++ dom delta') (gs ++ gs') = some sigma ->
      ps.length = vs.length ->
      Step d (.superCall v0 (.cls c gs) gs' m vs)
        ((e0.tsubst sigma).vsubst ((thisName, v0) :: (ps.map Prod.snd).zip vs))
  | innerFound {c fs c' gs vs0 ms gs' m vs mt delta' t0 ps e0 sigma} :
      (Expr.newLocal c fs c' gs vs0 ms).isValue = true -> Ty.allGround gs' = true ->
      Expr.allValues vs = true ->
      List.Mem mt ms ->
      FreshM (avoidNames d (.call (.newLocal c fs c' gs vs0 ms) gs' m vs)) mt (.mk delta' t0 m ps e0) ->
      mkSubst (dom delta') gs' = some sigma ->
      ps.length = vs.length ->
      Step d (.call (.newLocal c fs c' gs vs0 ms) gs' m vs)
        ((e0.tsubst sigma).vsubst ((thisName, .newLocal c fs c' gs vs0 ms) :: (ps.map Prod.snd).zip vs))
  | innerSuper {c fs c' gs vs0 ms gs' m vs} :
      (Expr.newLocal c fs c' gs vs0 ms).isValue = true -> Ty.allGround gs' = true ->
      Expr.allValues vs = true ->
      (ms.map Method.name).contains m = false ->
      Step d (.call (.newLocal c fs c' gs vs0 ms) gs' m vs)
        (.superCall (.newLocal c fs c' gs vs0 ms) (.cls c' gs) gs' m vs)
  | searchSuper {v c gs gs' m vs sigma} {dec : Dec} :
      List.Mem dec d -> dec.name = c ->
      v.isValue = true -> Ty.allGround gs = true -> Ty.allGround gs' = true -> Expr.allValues vs = true ->
      (dec.methods.map Method.name).contains m = false ->
      Not (c = objectName) ->
      mkSubst (dom dec.bounds) gs = some sigma ->
      Step d (.superCall v (.cls c gs) gs' m vs)
        (.superCall v (.cls dec.sup (Ty.substList sigma dec.supArgs)) gs' m vs)
  | ctx {E : ECtx} {e e'} :
      E.ok = true -> Step d e e' -> Step d (E.plug e) (E.plug e')
\end{lstlisting}
\paragraph*{Differences.}
Rule \textsc{simpleFieldFound} finds $f_i$ and $v_i$ by looking up $f$ in the pairs of the field names and the values $v_1,\ldots,v_n$; by A2 the field names are distinct.
Rule \textsc{simpleFound} requires as many parameters as arguments; the paper writes both as $x_1,\ldots,x_n$ and $v_1,\ldots,v_n$.
The $\mathsf{fresh}$ premises avoid fewer names (\texttt{avoidNames}, see the renaming above).
Rule \textsc{ifElse} writes $\emptyset\not\vdash\mathsf{typeOf}(v)<:\gen{C'}{\_}$ as: no type arguments make the judgement hold.
This is fine.

\subsection{Subtyping}
\paragraph*{Original.}
Rules \textsc{sub-refl}, \textsc{sub-inherit}, \textsc{sub-inherit-local} and \textsc{sub-var} of Section~\ref{sec:kinding}.
\paragraph*{Lean.}
\begin{lstlisting}[language=lean,basicstyle=\scriptsize\ttfamily]
inductive Sub (d : Decls) (delta : Bounds) : Ty -> Ty -> Prop
  | refl {t} : Sub d delta t t
  | inherit {c ts0 t' sigma} {dec : Dec} :
      List.Mem dec d -> dec.name = c ->
      mkSubst (dom dec.bounds) ts0 = some sigma ->
      Sub d delta ((Ty.cls dec.sup dec.supArgs).subst sigma) t' ->
      Sub d delta (.cls c ts0) t'
  | inheritLocal {c fs c' ts' es ms ts t' sigma} :
      Declared d (.loc c fs c' ts' es ms) ->
      mkSubstF fs ts = some sigma ->
      Sub d delta ((Ty.cls c' ts').subst sigma) t' ->
      Sub d delta (.cls c ts) t'
  | var {x t t'} :
      delta.lookup x = some t -> Sub d delta t t' -> Sub d delta (.var x) t'
\end{lstlisting}
\paragraph*{Differences.}
None beyond Section~\ref{app:lean:common}.

\subsection{Kinding}
\paragraph*{Original.}
Rules \textsc{wf-var} and \textsc{wf-class} of Section~\ref{sec:kinding}.
\paragraph*{Lean.}
\begin{lstlisting}[language=lean,basicstyle=\scriptsize\ttfamily]
inductive Kind (d : Decls) (delta : Bounds) : Ty -> Prop
  | var {x} : (delta.lookup x).isSome = true -> Kind d delta (.var x)
  | cls {c ts delta0 ts0 sigma} :
      scope d c ts = some (delta0, ts0) ->
      (forall t, List.Mem t ts -> Kind d delta t) ->
      mkSubst (dom delta0) ts0 = some sigma ->
      (forall p, List.Mem p (ts0.zip (Ty.substList sigma (cod delta0))) -> Sub d delta p.1 p.2) ->
      Kind d delta (.cls c ts)
\end{lstlisting}
\paragraph*{Differences.}
None beyond Section~\ref{app:lean:common}.

\subsection{Checked Cast}
\paragraph*{Original.}
Rules \textsc{cast-up}, \textsc{cast-var} and \textsc{cast-down} of Section~\ref{sec:typing}.
\paragraph*{Lean.}
\begin{lstlisting}[language=lean,basicstyle=\scriptsize\ttfamily]
inductive Cast (d : Decls) (delta : Bounds) : Ty -> Ty -> Prop
  | up {t t'} : Sub d delta t t' -> Cast d delta t t'
  | var {x t t'} : delta.lookup x = some t -> Cast d delta t t' -> Cast d delta (.var x) t'
  | down {c ts c' ts0 xs ts'} :
      Sub d delta (.cls c ts) (.cls c' ts0) ->
      generic d c = some xs ->
      Sub d [] (.cls c (vars xs)) (.cls c' ts') ->
      (forall x, List.Mem x xs -> List.Mem x (Ty.ftvList ts')) ->
      Cast d delta (.cls c' ts0) (.cls c ts)
\end{lstlisting}
\paragraph*{Differences.}
None beyond Section~\ref{app:lean:common}.

\subsection{$\mathsf{overrideOk}$ and $\mathsf{consistent}$}
\paragraph*{Original.}
The definitions of $\mathsf{overrideOk}(\Delta,T,\pl{M})$, $\mathsf{overrideOk}(T,M)$ (Section~\ref{sec:classtyping}) and $K\gap\mathsf{consistent}$ (Section~\ref{sec:typing}).
\paragraph*{Lean.}
\begin{lstlisting}[language=lean,basicstyle=\scriptsize\ttfamily]
def OverrideOk1 (d : Decls) (t : Ty) (mt : Method) : Prop :=
  match t, mt with
  | .cls c ts, .mk delta' t0' m ps' _ =>
    (exists dec : Dec, List.Mem dec d /\ dec.name = c /\ Not (List.Mem m (dec.methods.map Method.name)))
    \/ (exists (dec : Dec) (delta'' : Bounds) (t0 : Ty) (ps : List (Prod Ty Name)) (e : Expr) (sigma : TSubst),
          List.Mem dec d /\ dec.name = c /\
          List.Mem (Method.mk delta'' t0 m ps e) dec.methods /\
          mkSubst (dom dec.bounds ++ dom delta'') (ts ++ vars (dom delta')) = some sigma /\
          Ty.substList sigma (cod delta'') = cod delta' /\
          Ty.substList sigma (t0 :: ps.map Prod.fst) = t0' :: ps'.map Prod.fst)
    \/ (exists fs c0 ts0 es ms delta'' t0 ps e sigmaF sigma,
          Declared d (.loc c fs c0 ts0 es ms) /\
          List.Mem (Method.mk delta'' t0 m ps e) ms /\
          mkSubstF fs ts = some sigmaF /\
          mkSubst (dom delta'') (vars (dom delta')) = some sigma /\
          Ty.substList (sigmaF ++ sigma) (cod delta'') = cod delta' /\
          Ty.substList (sigmaF ++ sigma) (t0 :: ps.map Prod.fst) = t0' :: ps'.map Prod.fst)
  | .var _, _ => False
def OverrideOk (d : Decls) (delta : Bounds) (t : Ty) (ms : List Method) : Prop :=
  forall mt c ts, List.Mem mt ms -> Sub d delta t (.cls c ts) -> OverrideOk1 d (.cls c ts) mt
def Consistent (d : Decls) : K -> Prop
  | .loc c fs c0 ts _ ms =>
    exists fs' ts' es' ms' sigma,
      Declared d (.loc c fs' c0 ts' es' ms') /\
      mkSubstF fs' (fs.map Funnel.abs) = some sigma /\
      ts = Ty.substList sigma ts' /\
      (forall m, List.Mem m (ms'.map Method.name) <-> List.Mem m (ms.map Method.name)) /\
      forall mt, List.Mem mt ms -> OverrideOk1 d (.cls c (fs.map Funnel.abs)) mt
  | .top _ => False
\end{lstlisting}
\paragraph*{Differences.}
The three cases of $\mathsf{overrideOk}(T,M)$ are the three disjuncts.
The comparison of the types $T_0,\ldots,T_n$ with $T'_0,\ldots,T'_n$ as lists also requires the same number of parameters, as the paper does by writing $n$ in both methods.
In the last case, the substitution $[\pl{F},\Delta''=\pl{T},\mathrm{dom}(\Delta')]$ is the concatenation of the substitution of the funnels and the substitution of $\Delta''$; the domains are disjoint by A7.
$\mathrm{dom}(\pl{M}')=\mathrm{dom}(\pl{M})$ is an equivalence of membership.
This is fine.

\subsection{Typing}
\paragraph*{Original.}
The rules of Section~\ref{sec:typing}, and rule \textsc{new-local-java} of Section~\ref{sec:java}.
\paragraph*{Lean.}
\begin{lstlisting}[language=lean,basicstyle=\scriptsize\ttfamily]
mutual
inductive Typing (d : Decls) (java : Bool) : Bounds -> Ctx -> Expr -> Ty -> Prop
  | subs {delta gamma e t1 t2} :
      Typing d java delta gamma e t1 -> Kind d delta t2 -> Sub d delta t1 t2 ->
      Typing d java delta gamma e t2
  | var {delta gamma x t} :
      gamma.lookup x = some t -> Typing d java delta gamma (.var x) t
  | newObject {delta gamma} :
      Typing d java delta gamma (.new objectName [] []) (.cls objectName [])
  | declField {delta gamma e c ts' t f sigma} {dec : Dec} :
      Typing d java delta gamma e (.cls c ts') ->
      List.Mem dec d -> dec.name = c ->
      List.Mem (t, f) dec.fields ->
      mkSubst (dom dec.bounds) ts' = some sigma ->
      Typing d java delta gamma (.field e f) (t.subst sigma)
  | call {delta gamma e0 c0 ts0 delta0 ts2 delta1 t0 m ps body ts1 es sigma} :
      Typing d java delta gamma e0 (.cls c0 ts0) ->
      scope d c0 ts0 = some (delta0, ts2) ->
      MethodIn d (.mk delta1 t0 m ps body) c0 ->
      mkSubst (dom delta0 ++ dom delta1) (ts2 ++ ts1) = some sigma ->
      (forall t, List.Mem t ts1 -> Kind d delta t) ->
      (forall p, List.Mem p (ts1.zip (Ty.substList sigma (cod delta1))) -> Sub d delta p.1 p.2) ->
      es.length = ps.length ->
      (forall p, List.Mem p (es.zip (Ty.substList sigma (ps.map Prod.fst))) ->
        Typing d java delta gamma p.1 p.2) ->
      Typing d java delta gamma (.call e0 ts1 m es) (t0.subst sigma)
  | superCall {delta gamma e0 c ts c0 ts0 delta1 t0 m ps body ts1 es sigma} {dec0 dec : Dec} :
      Sub d delta (.cls c ts) (.cls c0 ts0) ->
      List.Mem dec0 d -> dec0.name = c0 ->
      MethodIn d (.mk delta1 t0 m ps body) c0 ->
      mkSubst (dom dec0.bounds ++ dom delta1) (ts0 ++ ts1) = some sigma ->
      List.Mem dec d -> dec.name = c ->
      Kind d delta (.cls c ts) ->
      (forall t, List.Mem t ts1 -> Kind d delta t) ->
      (forall p, List.Mem p (ts1.zip (Ty.substList sigma (cod delta1))) -> Sub d delta p.1 p.2) ->
      Typing d java delta gamma e0 (.cls c ts) ->
      es.length = ps.length ->
      (forall p, List.Mem p (es.zip (Ty.substList sigma (ps.map Prod.fst))) ->
        Typing d java delta gamma p.1 p.2) ->
      Typing d java delta gamma (.superCall e0 (.cls c ts) ts1 m es) (t0.subst sigma)
  | newTop {delta gamma c ts' es1 es2 sigma t''} {dec : Dec} :
      List.Mem dec d -> dec.name = c ->
      Kind d delta (.cls c ts') ->
      mkSubst (dom dec.bounds) ts' = some sigma ->
      es1.length = dec.fields.length ->
      (forall p, List.Mem p (es1.zip (Ty.substList sigma (dec.fields.map Prod.fst))) ->
        Typing d java delta gamma p.1 p.2) ->
      Typing d java delta gamma (.new dec.sup (Ty.substList sigma dec.supArgs) es2) t'' ->
      Typing d java delta gamma (.new c ts' (es1 ++ es2)) (.cls c ts')
  | newLocal {delta gamma c fs c0 ts es ms delta' ts'' t''} :
      java = false ->
      Consistent d (.loc c fs c0 ts es ms) ->
      pairing delta fs (fs.map Funnel.abs) = some (delta', ts'') ->
      (forall x, List.Mem x (Ty.ftvList (cod delta')) -> List.Mem x (dom delta')) ->
      Kind d delta' (.cls c (fs.map Funnel.abs)) ->
      Kind d delta' (.cls c0 ts) ->
      Typing d java delta gamma (.new c0 ts es) t'' ->
      (forall mt, List.Mem mt ms ->
        MethodOk d java delta'
          (restrict gamma (Ty.ftvList (fs.map Funnel.abs)) ++ [(thisName, .cls c (fs.map Funnel.abs))]) mt) ->
      OverrideOk d delta' (.cls c0 ts) ms ->
      Typing d java delta gamma (.newLocal c fs c0 ts es ms) (.cls c (fs.map Funnel.abs))
  | newLocalJava {delta gamma c c0 ts es ms t''} :
      java = true ->
      Kind d delta (.cls c0 ts) ->
      Typing d java delta gamma (.new c0 ts es) t'' ->
      (forall mt, List.Mem mt ms ->
        MethodOk d java delta (withoutThis gamma ++ [(thisName, .cls c [])]) mt) ->
      OverrideOk d delta (.cls c0 ts) ms ->
      Typing d java delta gamma (.newLocal c [] c0 ts es ms) (.cls c [])
  | ifInst {delta gamma e0 t0 t x e1 e2 t'} :
      Typing d java delta gamma e0 t0 -> Kind d delta t -> Cast d delta t0 t ->
      Typing d java delta (gamma ++ [(x, t)]) e1 t' -> Typing d java delta gamma e2 t' ->
      Typing d java delta gamma (.ifInst e0 t x e1 e2) t'
inductive MethodOk (d : Decls) (java : Bool) : Bounds -> Ctx -> Method -> Prop
  | mk {delta gamma delta' t0 m ps e} :
      (forall t, List.Mem t (cod delta') -> Kind d (delta ++ delta') t) ->
      (forall t, List.Mem t (t0 :: ps.map Prod.fst) -> Kind d (delta ++ delta') t) ->
      Typing d java (delta ++ delta') (gamma ++ ps.map fun p => (p.2, p.1)) e t0 ->
      MethodOk d java delta gamma (.mk delta' t0 m ps e)
end
\end{lstlisting}
\texttt{restrict gamma xs} is $\Gamma|_{\pl{X}}\setminus\this$ and \texttt{withoutThis gamma} is $\Gamma\setminus\this$.
\paragraph*{Differences.}
One relation with the flag \texttt{java} models both variants of Section~\ref{sec:java}: rule \textsc{new-local} requires \texttt{java = false}, and rule \textsc{new-local-java} requires \texttt{java = true}.
$\zeta(\mathsf{funnelled})$ is \texttt{Typing d false} and $\zeta(\mathsf{Java})$ is \texttt{Typing d true}.
Rule \textsc{call} substitutes $T_0,\ldots,T_n$ directly, instead of naming $T'_0,\ldots,T'_n$.
Rule \textsc{new-top} splits the arguments into the arguments for the fields, \texttt{es1}, and the arguments for the superclass, \texttt{es2}.
This is fine.

\subsection{Typing of Programs, Classes and Methods}
\paragraph*{Original.}
Rules \textsc{prog-ok}, \textsc{method-ok} and \textsc{class-ok} of Section~\ref{sec:classtyping}.
\paragraph*{Lean.}
\texttt{MethodOk} is in the listing of the typing above.
\begin{lstlisting}[language=lean,basicstyle=\scriptsize\ttfamily]
def ClassOk (d : Decls) (java : Bool) (dec : Dec) : Prop :=
  (forall t, List.Mem t (cod dec.bounds) -> Kind d dec.bounds t) /\
  (forall t, List.Mem t (Ty.cls dec.sup dec.supArgs :: dec.fields.map Prod.fst) -> Kind d dec.bounds t) /\
  (forall mt, List.Mem mt dec.methods ->
    MethodOk d java dec.bounds [(thisName, .cls dec.name (vars (dom dec.bounds)))] mt) /\
  OverrideOk d dec.bounds (.cls dec.sup dec.supArgs) dec.methods
def ProgOk (d : Decls) (java : Bool) (e : Expr) : Prop :=
  (forall dec, List.Mem dec d -> ClassOk d java dec) /\ exists t, Typing d java [] [] e t
\end{lstlisting}
\paragraph*{Differences.}
\texttt{MethodOk} is mutual with \texttt{Typing}, since rule \textsc{new-local} has a premise $\mok{\Delta'}{\Gamma'}{\pl{M}}$.
\texttt{ClassOk} and \texttt{ProgOk} are definitions, since they have one rule each.
This is fine.

\subsection{Statements and High Level Proofs}
\label{app:lean:statements}
\begin{lstlisting}[language=lean,basicstyle=\scriptsize\ttfamily]
def DeclsOk (d : Decls) : Prop := wfDecls d = true /\ forall dec, List.Mem dec d -> ClassOk d false dec
theorem stepWellFormed {d e e' t} (hd : DeclsOk d) (he : wfExpr e = true)
    (ht : Typing d false [] [] e t) (hs : Step d e e') : wfExpr e' = true
theorem progress {d e t} (hd : DeclsOk d) (he : wfExpr e = true)
    (ht : Typing d false [] [] e t) :
    e.isValue = true \/ exists e', Step d e e'
theorem preservation {d e e' t} (hd : DeclsOk d) (he : wfExpr e = true)
    (ht : Typing d false [] [] e t) (hs : Step d e e') : Typing d false [] [] e' t
theorem typeSoundness {d e0 e t} (hd : DeclsOk d) (he0 : wfExpr e0 = true)
    (ht : Typing d false [] [] e0 t) (hs : Steps d e0 e) :
    Typing d false [] [] e t /\ (e.isValue = true \/ exists e', Step d e e')
\end{lstlisting}
\texttt{progress}, \texttt{preservation} and \texttt{typeSoundness} are Progress, Preservation and Type Soundness of the section on the (un)soundness of method local classes.
\texttt{DeclsOk d} is the global \decls such that $\mathsf{Dec}\gap\mathrm{Ok}$ for every $\mathsf{Dec}\in\decls$, together with A1-A11 on \decls.
\texttt{stepWellFormed} has no counterpart in the paper (Section~\ref{app:lean:common}).
\texttt{typeSoundness} is proved by induction on \texttt{Steps}, from the other three theorems.
The planned proofs of the other three theorems are:
\begin{enumerate}
\item \textbf{Basic lemmas.} \texttt{Ty.beq} is equality; \texttt{lookup} on a substitution built by \texttt{mkSubst}; by A1, \texttt{declaredKs}, \texttt{scope} and \texttt{generic} have one entry per class name.
\item \textbf{Renaming.} \texttt{Expr.freshen} renames bound names consistently, so typing, kinding and well-formedness are invariant under it; the result binds only names of \texttt{ns}.
\item \textbf{Weakening} of $\Gamma$ and $\Delta$, for typing, kinding and subtyping.
\item \textbf{Substitution of type variables.} If $\Delta,\Delta'$ types $e$, and ground types $\pl{G}$ satisfy the bounds $\Delta'$, then $e[\Delta'=\pl{G}]$ is typed under $\Delta$ with the substituted type.
Kinding and subtyping are preserved, but the checked cast is not: in rule \textsc{cast-var}, after $X$ becomes $G$, the source $G$ may not reach the target through \textsc{cast-up} or \textsc{cast-down}.
For the typing of rule \textsc{if-instanceof} the lemma uses rule \textsc{subs} to lift $e_0$ from $G$ to the substituted bound of $X$, and applies \textsc{cast-down} from there.
\item \textbf{Substitution of variables.} Values are closed and typed under $\emptyset;\emptyset$, so substituting values for variables of the same types preserves typing; \Q@this@ is not substituted into the methods of a local class.
\item \textbf{Inversion} of typing for each form of expression, through rule \textsc{subs}.
\item \textbf{Canonical forms.} A value of type $\gen{C}{\pl{T}}$ has a class that is a subclass of $C$.
\item \textbf{Copies of a local class.} By $\mathsf{consistent}$, the methods of a copy have the signatures of the declaration with the funnelled type variables substituted; so rules \textsc{call} and \textsc{sub-inherit-local}, which read the declaration, agree with rule \textsc{innerFound}, which runs the copy.
\item \textbf{Checked cast.} If $\emptyset\vdash\mathsf{typeOf}(v)<:\gen{C}{\_}$, $v$ has type $T_0$ and $\cast{\emptyset}{T_0}{\gen{C}{\pl{T}}}$, then $v$ has type $\gen{C}{\pl{T}}$.
The proof uses that the chain of superclasses is unique and that $\pl{X}\subseteq\ftv{\pl{T}'}$ fixes $\pl{T}$.
\item \textbf{stepWellFormed} by induction on \texttt{Step}: the fresh names keep A7, and the other requirements are local.
\item \textbf{progress} by induction on typing, with canonical forms; a call terminates its search at a class declaring the method, before \Q@Object@, since typing found the method in a superclass.
\item \textbf{preservation} by induction on \texttt{Step}, with the lemmas above.
\end{enumerate}

\subsection{Validation Examples}
\label{app:lean:examples}
\leanfile{Examples.lean} defines the class tables below and the examples.
A Lean theorem marked \emph{proved} is accepted by Lean without \texttt{sorry}; a theorem marked \emph{stated} has \texttt{sorry}, and needs the inversion lemmas of Section~\ref{app:lean:statements}.
\begin{lstlisting}
class Object<> extends Object<>{}
class A<> extends Object<>{}
class B<> extends Object<>{ B<> onlyB(){ return new B<>(); } }
class Box<X extends Object<>> extends Object<>{ X val; X get(){ return this.val; } }
class C1<> extends Object<>{ A<> m(){ return new A<>(); } }
class C2<> extends C1<>{}                                         // d1: the classes above
class Foo<X extends Object<>> extends Object<>{ X x(){ return this.x(); } }
class Maker<> extends Object<>{ <Y extends Object<>> Foo<Y> mk(Y y){ return
  new MyFoo<Y funnel> extends Foo<Y>(){ Y x(){ return y; } }; } }  // d2: Object, A, B, Foo, Maker
class Broken<> extends Object<>{ <X extends Object<>> Foo<X> m(Object<> other, X x){ return
  if (other instanceof MyFoo<> w){ w }
  else { new MyFoo<> extends Foo<X>(){ X x(){ return x; } } }; } }  // d3: Object, A, B, Foo, Broken
class Leaker<> extends Object<>{
  <Y extends Object<>> MyFoo2<> mk2(Y y){ return new MyFoo2<> extends Foo<Y>(){ Y x(){ return y; } }; }
  <Y extends Object<>> Foo<Y> leak(MyFoo2<> z){ return z; } }     // d4: Object, A, B, Foo, Leaker
\end{lstlisting}
\begin{description}
\item[W1] \texttt{w1}, proved: d1 satisfies A1-A11.
\item[W2] \texttt{w2}, proved: a method with two parameters \Q@x@ violates A4.
\item[W3] \texttt{w3}, proved: an \Q@if@ binding the parameter \Q@x@ again violates A7.
\item[W4] \texttt{w4}, proved: an undeclared superclass violates A9 and A10.
\item[W5] \texttt{w5}, proved: d2 satisfies A1-A11.
\item[W6] \texttt{w6}, proved: d3 and the main expression of Example~1 satisfy A1-A11.
\item[W7] \texttt{w7}, proved: d4 satisfies A1-A11.
\item[R1] \texttt{r1}, proved: \Q@new Box<A<>>(new A<>()).val@ $\rightarrow$ \Q@new A<>()@ by \textsc{simpleFieldFound}.
\item[R2] \texttt{r2}, proved: \Q@new C2<>().m()@ $\rightarrow$ \Q@new C2<>().[C2<>]m()@ by \textsc{simpleAddSuper}.
\item[R3] \texttt{r3}, proved: \Q@new C2<>().[C2<>]m()@ $\rightarrow$ \Q@new C2<>().[C1<>]m()@ by \textsc{searchSuper}.
\item[R4] \texttt{r4}, proved: \Q@new C2<>().[C1<>]m()@ $\rightarrow$ \Q@new A<>()@ by \textsc{simpleFound}.
\item[R5] \texttt{r5}, proved: R1 inside the receiver of a call, by \textsc{ctx}.
\item[R6] \texttt{r6Then}, proved: \Q@if (new A<>() instanceof A<> x){ x } else { new B<>() }@ $\rightarrow$ \Q@new A<>()@ by \textsc{ifThen}.
\item[R6] \texttt{r6Else}, stated: the same with \Q@B<>@ as the test $\rightarrow$ \Q@new B<>()@ by \textsc{ifElse}.
\item[R7] \texttt{r7}, proved: \Q@new Maker<>().[Maker<>]<A<>>mk(new A<>())@ $\rightarrow$ \Q@new MyFoo<A<>> extends Foo<A<>>(){ A<> x(){ return new A<>(); } }@ by \textsc{simpleFound}, renaming \Q@Y@ and \Q@y@ into \Q@Y1@ and \Q@y1@.
\item[R8] \texttt{r8}, proved: the value of R7 \Q@.x()@ $\rightarrow$ \Q@new A<>()@ by \textsc{innerFound}.
\item[R9] \texttt{r9}, proved: the value of R7 \Q@.other()@ $\rightarrow$ \Q@.[Foo<A<>>]other()@ by \textsc{innerSuper}.
\item[R10] \texttt{r10}, proved: the value of R7 \Q@.f@ $\rightarrow$ \Q@new Foo<A<>>().f@ by \textsc{innerFieldSuper}.
\item[R11] \texttt{r11}, proved: R2, R3, R4 in sequence.
\item[T1] \texttt{t1}, proved: \Q@new Box<A<>>(new A<>()).val : A<>@.
\item[T2] \texttt{t2}, stated: \Q@new Box<A<>>(new B<>())@ is ill typed: \Q@B<>@ is not a subtype of \Q@A<>@.
\item[C1] \texttt{c1}, proved: $\cast{\QQ{X extends Object<>}}{\QQ{Foo<X>}}{\QQ{MyFoo<X>}}$ in d2.
\item[C2] \texttt{c2}, stated: not $\cast{\QQ{X extends Object<>}}{\QQ{Object<>}}{\QQ{MyFoo<X>}}$ in d2.
\item[M1] \texttt{m1}, stated: class \Q@Maker@ is Ok under funnelled typing.
\item[J1] \texttt{j1}, stated: Example~1 (d3) is well typed under Java typing.
\item[J2] \texttt{j2}, stated: Example~1 reduces to \Q@new A<>().onlyB()@, which is stuck.
\item[N1] \texttt{n1}, stated: Example~1 is ill typed under funnelled typing.
\item[L1] \texttt{l1}, stated: class \Q@Leaker@ is Ok under Java typing.
\item[L2] \texttt{l2}, stated: class \Q@Leaker@ is not Ok under funnelled typing.
\end{description}
\smallskip

\noindent In the paper, the verdict of each example follows from the rule named in the list, or, for W1-W7, from the requirements named in the list.
For R7, the paper renames with $\mathsf{fresh}$ the type variable \Q@Y@ and the parameter \Q@y@ of \Q@mk@, and then substitutes \Q@A<>@ and \Q@new A<>()@; the method \Q@x@ of the local class binds no names.
For T2, rule \textsc{new-top} needs \Q@new B<>() : A<>@, and no rule gives it.
For C2, rule \textsc{cast-down} needs $\QQ{X}\in\ftv{\emptyset}$, since \Q@MyFoo<Y>@ has the supertype \Q@Object<>@ with no type arguments.
For J1, J2 and N1, Section~\ref{sec:java} gives the derivation; N1 fails at $\kind{\emptyset}{\QQ{Foo<X>}}$ in rule \textsc{new-local}, since \Q@MyFoo@ funnels nothing.
For L1 and L2, Section~\ref{sec:java} discusses the method \Q@leak@; L2 fails at the local class in \Q@mk2@, as N1 does.
